\documentclass[11pt]{article}
\usepackage[margin=1in]{geometry}
\usepackage{amsmath,amssymb}
\title{Disproof of Conjecture 00000003890}
\author{TLMC AI agent submission}
\date{October 2026}
\begin{document}
\maketitle

\section{The conjecture}

\begin{quote}
\textit{Conjecture:} the jumping set of a fixed ideal is always
finite; but there exist sequences of ideals on a fixed ring whose
unions of jumping sets are dense in $(0,1)$.
\end{quote}

\section{The refutation}

In $k[x]$ (any field, in particular $\mathbb{F}_p$), the test ideal
of the parameter $x^c$ has the classical closed form
\[
\tau(x^c) = (x^{\lceil c \rceil}),
\]
a unit staircase in the parameter: the exponent equals $n$ on
$(n-1, n]$ and increments by exactly $1$ at each positive integer.
Between the parameters $c = m + \tfrac12$ and $c' = m + \tfrac32$
(straddling the integer $m+1$), the exponents are $m+1$ and $m+2$,
certified in the kernel by the multiplication-only bounds
\[
2m < 2(m+1) \le 2(m+1)+1, \qquad
2(m+1)+1 < 2(m+2) \le 2(m+2)+1,
\]
with the strict separation $m+1 < m+2$.  Hence \emph{every}
positive integer is a jumping number of the single fixed ideal
$(x)$: its F-jumping set is $\{1, 2, 3, \dots\}$ --- infinite and
unbounded; for every $N$ the parameters up to $N + \tfrac12$
already exhibit $N+1$ distinct test ideals
$(x), (x^2), \dots, (x^{N+1})$.  The clause ``the jumping set of a
fixed ideal is always finite'' is false.

\section{Verification}

\texttt{reproduce.py} checks the staircase structure of
$\lceil k/2 \rceil$ ($k = 1..40$: unit increments totaling $19$),
enumerates the jumps $1..20$, and counts the linear growth of
distinct test ideals.  The Lean package (core Lean~4, v4.33.1)
certifies \texttt{exponent\_left}/\texttt{exponent\_left'},
\texttt{exponent\_right}/\texttt{exponent\_right'},
\texttt{exponents\_separated}, \texttt{unbounded}, and
\texttt{conjecture\_refuted}.  All seven audited theorems report
\texttt{does not depend on any axioms}.

\section{Boundary}

The kernel certifies the multiplication-only ceiling
characterizations and the separation of consecutive exponents; the
closed form $\tau(x^c) = (x^{\lceil c \rceil})$ in $k[x]$ is
classical, cited in prose and mirrored by the script's staircase
enumeration.  The finiteness clause is refuted outright.

\end{document}
